\section{Likelihood truncation and selected dyadic cells}
\label{sec:l-likelihood-truncation}

This section gives a forward proof of the precision lower bound using a
cube chart of the sphere.  At a fixed incoming state, the likelihood
of a block-ending state at a signal is the probability that the fresh
block ends in that state when the signal is fixed.  A block may
concentrate the conditional signal law, but a Gaussian projection
estimate limits the increase in a typical state likelihood.  We then select a dyadic cell
for analysis and condition on the event that this cell was selected.
The cost of that conditioning and the depth of the cell enter one
density bound.  Iterating the bound prevents a small number of blocks
from localizing the signal to the accuracy required by the output.

The use of high Gaussian projection moments has a methodological
precedent in Sharan, Sidford, and Valiant
\cite[Section~7, Lemma~12]{SSV2019}.  We prove the exact-label estimate
used here, including its dependence on the mass of a restricted set.
The lower bound in this section starts from success for every signal:
that hypothesis supplies both the cube-prior average used by the
forward argument and the uniform-sphere average used at constant
accuracy.

\subsection{The learner, the cube prior, and measurable rules}
\label{subsec:l-lt-model}

We use the finite-state learner of Section~\ref{l:model}, including
its jointly measurable experiment, joint independence of the shared
rule seed and initialization from the signal and rows, and permitted
per-seed completed-measurability.  The deterministic horizon is
\(T\ge0\), and the output uses the terminal state, stopping index,
and fresh randomness.

\begin{theorem}[Precision lower bound from likelihood truncation]
\label{thm:l-lt-precision}
Let $M(d)=o(d^2)$ be a nonnegative integer-valued memory bound, and let
$0<\epsilon(d)\le1/10$.  Suppose a family of learners in the model
above satisfies
\[
 \mathbb P\!\left\{\arccos\langle\widehat s,s\rangle
              \le\epsilon(d)\right\}\ge\frac23
 \qquad\text{for every }s\in S^{d-1}.
\]
There is an absolute $c>0$ such that, for all sufficiently large $d$,
\[
 T(d)\ge c\,d\log_2\frac1{\epsilon(d)}.
\]
The eventual dimension threshold may depend on the memory sequence.
\end{theorem}

Put $p=d-1$ and consider the normalized cube law
\begin{equation}\label{l:lt:chart}
 Q_*=\left[-\frac1{2\sqrt p},\frac1{2\sqrt p}\right)^p,\qquad
 \mu=\operatorname{Unif}(Q_*),\qquad
 s(z)=\bigl(z,\sqrt{1-\|z\|^2}\bigr).
\end{equation}
The parameter law has density \(p^{p/2}\) on the cube.  This is the
cube chart of Lemma~\ref{ld:cube-geometry}, with \(n\) renamed \(p\).
That lemma proves that \(s\) is \(2\)-Lipschitz and, for every fixed
unit output \(v\),
\begin{equation}\label{l:lt:output-mass}
 \mu\{z:\arccos\langle v,s(z)\rangle\le\epsilon\}
 \le\bigl(\sqrt{2\pi e}\,\epsilon\bigr)^p\le(5\epsilon)^p.
\end{equation}
Write \(V_h\) for the volume of the unit ball in \(\mathbb R^h\).
The same Gaussian-integral proof gives
\(V_h\le(\sqrt{2\pi e}/\sqrt h)^h<(5/\sqrt h)^h\), which will also
be used below.

The following reduction puts the learner in the form required by the
block argument.  It makes the order of the two operations explicit:
first fix the shared rule-selection randomness, then choose Borel
versions for the resulting cube-prior experiment.

\begin{lemma}[A deterministic Borel learner for the cube experiment]
\label{lem:l-lt-borel}
Fix $d\ge2$, $M$, and a finite horizon $T$.  If a learner in the stated
model has average success at least $2/3$ under \eqref{l:lt:chart}, then
there is a deterministic learner with Borel transition and stopping
rules, the same state and horizon bounds, and a fixed unit output at
each terminal state/index label, whose average success under that prior
is at least $3/5$.
\end{lemma}
\begin{proof}
Choose a shared-seed value with conditional cube-prior success greater
than \(3/5\).  The joint experiment in Section~\ref{l:model} makes this
conditioning meaningful, and the conditional initial-state law remains
independent of the signal and rows.  The derivative of the chart has
the identity in its first \(p\) rows, so Lemma~\ref{l:prior-ac} proves
absolute continuity of its one-pair law.  Lemma~\ref{l:borel} therefore
provides everywhere Borel finite kernels preserving this fixed-seed
cube-prior joint law of the signal, path, stopping index, and output.

Use the finite-randomness reduction following that lemma.  Sample the
conditional initial state, independent uniforms for all finite choices
indexed by state and index (including an input-free index-zero stopping
choice), and one output from each terminal state/index law.  These
finitely many choices are independent of the inputs.  Fixing a
realization by averaging retains cube-prior success at least \(3/5\).
The resulting state rules are Borel, and the state and horizon bounds
are unchanged.
\end{proof}

\subsection{The projection estimate inside a regular cell}
\label{subsec:l-lt-projection}

We next control one block, starting from a parameter law with a bound
on concentration at every smaller dyadic scale.  For sufficiently
large $d$, set
\begin{equation}\label{l:lt:parameters}
 n=\lfloor d/32\rfloor,\qquad
 k=\lfloor d/8\rfloor,\qquad
 \lambda=p/2,\qquad Q=\lambda-k,\qquad \alpha=\frac nQ.
\end{equation}
Here $n$ is the number of rows in a block and $k$ is the moment order.
For example, $d\ge64$ ensures $n\ge1$, $k\ge2$,
$Q\ge d/3$, and $\alpha\le1/4$.

Subdivide $Q_*$ dyadically in every coordinate, using half-open cells
so that every parameter belongs to one cell at each depth.  A cell
$C$ at depth $l$ has side length $R/\sqrt p$, where $R=2^{-l}$.
A descendant at relative depth $j$ is a depth-$l+j$ cell inside $C$.
Call a Borel probability measure $P$ supported in $C$ regular when
\begin{equation}\label{l:lt:regular}
 P(V)\le2^d2^{-\lambda j}
 \quad\text{for every descendant $V$ at relative depth $j\ge0$}.
\end{equation}

We need a projection estimate that retains the probability of a
restricted parameter set.  In the likelihood cutoff, this set will
consist of signals for which a particular state is unusually likely.
The factor $\rho^{1-\alpha}$ below lets the small mass $\rho$ of
that set survive both the projection and the choice of a state.

\begin{lemma}[Restricted mass and exact projection moments]
\label{lem:l-lt-projection}
There are absolute constants $C_0,C_1,C_2\ge1$ such that the following
holds with \eqref{l:lt:parameters} for all sufficiently large $d$.
Let $P$ be regular in a cell $C$ of scale $R$.  For every
$x\in\mathbb R^d$, $u>0$, and affine subspace $H\subseteq\mathbb R^d$
of dimension at most $k$,
\begin{align}
 P\{z:s(z)\in B(x,u)\}
 &\le C_0^d(u/R)^\lambda, \label{l:lt:ball}\\
 P\{z:\operatorname{dist}(s(z),H)\le u\}
 &\le(C_1u/R)^Q. \label{l:lt:tube}
\end{align}
For a Borel set $E\subseteq C$ of mass $\rho=P(E)>0$,
\begin{equation}\label{l:lt:inverse}
 \int_E\operatorname{dist}(s(z),H)^{-n}\,dP(z)
 \le C_2^dR^{-n}\rho^{1-\alpha}.
\end{equation}

Let $\gamma_{n,d}$ denote standard Gaussian probability on
$\mathbb R^{n\times d}$ and put
$\beta(dA,dy)=\gamma_{n,d}(dA)\,dy$.  For $\eta>0$, let
$\varphi_\eta$ be the centered Gaussian density on $\mathbb R^n$ with
covariance $\eta^2I_n$, and define
\[
 J_\eta(A,y)=\int_E\varphi_\eta(y-As(z))\,dP(z).
\]
Then
\begin{equation}\label{l:lt:moment}
 \int J_\eta^k\,d\beta
 \le \bigl(C_2^dR^{-n}\bigr)^{k-1}
       \rho^{\,1+(k-1)(1-\alpha)}.
\end{equation}
For every Borel joint set $B\subseteq
\mathbb R^{n\times d}\times\mathbb R^n$ with $\beta(B)<\infty$,
the exact observations satisfy
\begin{equation}\label{l:lt:exact-set}
 \mathbb P_{z\sim P,\,A\sim\gamma_{n,d}}
       \{z\in E,\ (A,As(z))\in B\}
 \le \beta(B)^{1-1/k}
       \bigl(C_2^dR^{-n}\bigr)^{1-1/k}\rho^{1-\alpha}.
\end{equation}
The same assertions involving $E$ are trivial when $\rho=0$.
\end{lemma}
\begin{proof}
For $0<u<R$, choose an integer $j\ge1$ with
$R2^{-j}\le u<2R2^{-j}$, and write $\delta=R2^{-j}$.
Project $B(x,u)$ onto the first $p$ coordinates.  Each relative-depth
$j$ cell meeting the projected ball has diameter $\delta$, so the
union of those cells lies in a ball of radius $u+\delta<3\delta$.
The cells have disjoint interiors and volume
$(\delta/\sqrt p)^p$.  The unit-ball estimate used in
\eqref{l:lt:output-mass} bounds their number by
\[
 \frac{V_p(3\delta)^p}{(\delta/\sqrt p)^p}\le15^p.
\]
Using \eqref{l:lt:regular} for each cell gives
$P\{s(z)\in B(x,u)\}\le15^p2^d2^{-\lambda j}$.
Since $2^{-j}\le u/R$, this proves \eqref{l:lt:ball} after choosing
$C_0$ absolutely.  For $u\ge R$ the bound follows from $P(C)=1$.

Let $s_0$ be the lift of the center of $C$.  The cell has Euclidean
radius $R/2$ in parameter space, so the Lipschitz bound implies
$\|s(z)-s_0\|\le R\le2R$ for $z\in C$.  Orthogonal projection onto
$H$ is nonexpansive.  Thus the projections of all these lifted points
lie in the radius-$2R$ ball in $H$ about the projection of $s_0$.
Disjoint ball packing covers that ball by at most
$(1+4R/u)^k$ balls of radius $u$.  A point at distance at most $u$
from $H$ whose projection lies in one of these balls is within $2u$
of its center.  For $u\le R$, applying \eqref{l:lt:ball} to the cover
gives
\[
 P\{\operatorname{dist}(s(z),H)\le u\}
 \le (1+4R/u)^k C_0^d(2u/R)^\lambda
 \le C^d(u/R)^{\lambda-k}.
\]
Here and below $C$ denotes an absolute constant that may increase.
The identity
$Q=(d-1)/2-\lfloor d/8\rfloor\ge d/3$ for large $d$ lets us absorb
$C^d$ into $C_1^Q$.  For $u>R$, enlarge $C_1$ so that the right side
is at least one.  This proves \eqref{l:lt:tube}.  In particular every
affine subspace of the stated dimension has $P$-mass zero, by letting
$u$ decrease to zero.

To prove \eqref{l:lt:inverse}, apply the layer-cake formula to
$(\operatorname{dist}(s(z),H)/R)^{-n}\mathbf 1_E$.  At height $t>0$
its tail has mass at most
$\min\{\rho,C_1^Qt^{-Q/n}\}$.  With
$t_0=C_1^n\rho^{-n/Q}$,
\begin{align*}
 \int_0^\infty\min\{\rho,C_1^Qt^{-Q/n}\}\,dt
 &\le \rho t_0+C_1^Q\int_{t_0}^\infty t^{-Q/n}\,dt\\
 &=\frac{C_1^n}{1-\alpha}\rho^{1-\alpha}.
\end{align*}
Because $\alpha\le1/4$ and $n\le d$, an absolute $C_2$ bounds the
prefactor by $C_2^d$.  Restoring $R^{-n}$ proves the assertion.

We turn from affine distances to the projection moment.  Tonelli's
theorem expands $\int J_\eta^k\,d\beta$ over
$z_1,\ldots,z_k\in E$.  Write $s_i=s(z_i)$.  For fixed $A$ and this
tuple, integration of
$\prod_{i=1}^k\varphi_\eta(y-As_i)$ over the common label $y$ is the
density at zero of the $k-1$ differences from the first of the
independently mollified labels $As_i+\eta Z_i$.  After averaging over
$A$, the differences are centered Gaussian, independently across
rows, with per-row covariance
\begin{equation}\label{l:lt:covariance}
 G+\eta^2(I+\mathbf1\mathbf1^{\mathsf T}),\qquad
 G=\bigl(\langle s_i-s_1,s_j-s_1\rangle\bigr)_{i,j=2}^k.
\end{equation}
Their density at zero is
\[
 (2\pi)^{-n(k-1)/2}
 \det\!\left(G+\eta^2(I+\mathbf1\mathbf1^{\mathsf T})\right)^{-n/2}
 \le \det(G)^{-n/2},
\]
where the upper bound is interpreted as infinity when $G$ is
singular.  For nonsingular $G$, the inequality follows by factoring
$G^{1/2}$ and observing that all eigenvalues of
$I+\eta^2G^{-1/2}(I+\mathbf1\mathbf1^{\mathsf T})G^{-1/2}$
are at least one.  Gram--Schmidt gives
\[
 \det(G)^{-n/2}
 =\prod_{i=2}^k
   \operatorname{dist}\bigl(s_i,
      \operatorname{aff}(s_1,\ldots,s_{i-1})\bigr)^{-n}.
\]
The affine spaces here have dimension at most $k-2$.  The null-mass
consequence of \eqref{l:lt:tube}, applied successively, shows that
degenerate tuples have product $P$-measure zero.  Integrate the last
display from $z_k$ down to $z_2$ using \eqref{l:lt:inverse}, then
integrate $z_1$ with mass $\rho$.  The result is
\[
 \rho\bigl(C_2^dR^{-n}\rho^{1-\alpha}\bigr)^{k-1},
\]
which is \eqref{l:lt:moment}.

H\"older's inequality now bounds the mollified measure of $B$ by
\[
 \beta(B)^{1-1/k}
 \bigl(C_2^dR^{-n}\bigr)^{1-1/k}
 \rho^{\,1-\alpha+\alpha/k}.
\]
Since $\rho\le1$, this implies the weaker power of $\rho$ in
\eqref{l:lt:exact-set}.  For an open $B$, couple the mollified labels
as $As(z)+\eta Z$ using one independent standard Gaussian $Z$.
As $\eta$ decreases to zero, membership of $(A,As(z))$ in $B$
implies eventual membership of the approximants.  Fatou's lemma
therefore passes the same bound to the exact observation measure.
The Gaussian--Lebesgue measure $\beta$ is a locally finite regular
Borel measure on the Euclidean joint space.  Its outer regularity
extends the bound from open sets to every Borel $B$ of finite measure.
In particular the exact measure vanishes on $\beta$-null sets, so the
same bound extends to the completion of $\beta$ as well.  The mollifier
was used only in this calculation; the observation in
\eqref{l:lt:exact-set} is the exact label.
\end{proof}

We now apply the exact-set estimate to the signals for which one
state's likelihood exceeds a multiple of its marginal probability.
The proof restricts labels to a ball about the projected cell center.
The resulting sets of observations producing different states are
disjoint, and their total reference volume will control the sum over states.

\begin{lemma}[A regular-cell likelihood cutoff]
\label{lem:l-lt-likelihood}
There are absolute $a,c_1>0$ and a dimension threshold such that the
following holds.  Let $P$ be regular in a cell $C$, let
$A\sim\gamma_{n,d}$ be independent of $z\sim P$, and let
\[
 W=f(A,As(z))
\]
for a Borel map $f$ with at most $N$ values.  Suppose
$N^{1/k}\le2^d$, and set
\[
 q_w(z)=\mathbb P_A\{f(A,As(z))=w\},\qquad
 b_w=\int q_w\,dP.
\]
Then
\begin{equation}\label{l:lt:likelihood-tail}
 \mathbb P\{q_W(z)>2^{ad}b_W\}\le2^{-c_1d}.
\end{equation}
\end{lemma}
\begin{proof}
The functions $q_w$ are Borel by integration of the Borel routing
indicators.  For $b_w>0$ define
$E_w=\{z:q_w(z)>2^{ad}b_w\}$.  Markov's inequality gives
$P(E_w)\le2^{-ad}$.  A state with $b_w=0$ is reached with probability
zero and contributes nothing.

Retain the lift $s_0$ of the center of $C$ used in the preceding
proof.  Conditional on $z$, the vector $A(s(z)-s_0)$ is distributed
as $\|s(z)-s_0\|Z$ for a standard Gaussian $Z\in\mathbb R^n$.
Because $\|s(z)-s_0\|\le2R$,
\begin{equation}\label{l:lt:label-tail}
 \mathbb P\{\|As(z)-As_0\|>4R\sqrt n\}
 \le \mathbb P\{\|Z\|^2>4n\}
 \le 2^{n/2}e^{-n}.
\end{equation}
The last inequality uses
$\mathbb E e^{\|Z\|^2/4}=2^{n/2}$.  Since $n\ge d/64$ for large
$d$, the right side is exponentially small in $d$.

Define the joint domain
\[
 \mathcal D=\{(A,y):\|y-As_0\|\le4R\sqrt n\}
\]
and let $B_w=\mathcal D\cap f^{-1}(\{w\})$.  These Borel sets are
disjoint, and the volume bound gives
\[
 \sum_w\beta(B_w)\le\beta(\mathcal D)
 =V_n(4R\sqrt n)^n\le(20R)^n.
\]
Apply \eqref{l:lt:exact-set} to $E_w$ and $B_w$, and sum over $w$.
H\"older's inequality on the finite state set yields
\begin{align}
 \sum_w\mathbb P\{z\in E_w,(A,As(z))\in B_w\}
 &\le \bigl(C_2^dR^{-n}\bigr)^{1-1/k}
       2^{-ad(1-\alpha)}
       N^{1/k}\left(\sum_w\beta(B_w)\right)^{1-1/k}\notag\\
 &\le C_4^dN^{1/k}2^{-ad(1-\alpha)}
 \label{l:lt:width-tail}
\end{align}
for an absolute $C_4\ge1$.  The powers of $R$ cancel exactly.
Since $1-\alpha\ge3/4$ and $N^{1/k}\le2^d$, choose $a$ so that
$3a/4\ge\log_2 C_4+3$.  The last display is then at most $2^{-2d}$.
Adding \eqref{l:lt:label-tail} and reducing an absolute $c_1>0$ if
necessary proves \eqref{l:lt:likelihood-tail} for large $d$.
\end{proof}

\begin{remark}
The projection estimate also verifies the cube-prior absolute
continuity used in Lemma~\ref{lem:l-lt-borel}.  At the root,
$\mu(V)=2^{-pj}$ for a depth-$j$ cell, so $\mu$ is regular.
Taking $P=\mu$ and $E=Q_*$ in \eqref{l:lt:exact-set} shows that the
exact joint block law is absolutely continuous with respect to
$\beta$.  For the one-pair reference measure
$\lambda_d(dx,dy)=\gamma_d(dx)\,dy$ from Section~\ref{l:model},
the inverse image of a $\lambda_d$-null set under one-row projection
is $\beta$-null by the product structure and $\sigma$-finiteness.
The one-pair marginal is therefore absolutely continuous as well.
\end{remark}

\subsection{Restoring regularity by a selection event}
\label{subsec:l-lt-selection}

Fix a deterministic Borel learner, padded to $m$ blocks of $n$ rows.
At each block boundary let its state alphabet have size at most $N$,
with $N^{1/k}\le2^d$.  A halted path keeps a tag recording its terminal
state and index; the quantitative application below supplies this
padding.  For a fixed starting state, the next boundary state is a
Borel function of the whole exact block $(A,As(z))$.  Trailing rows
after a halt are generated and ignored.

For the proof only, a history records the boundary states, selected
dyadic cells, and survival of the deletions defined below.  At time
$t$ it is an event determined by the parameter and the first $t$
blocks.  We condition only on histories of positive probability.
Future rows remain independent standard Gaussian rows conditional on
such a history and the parameter, since those rows are independent of
the parameter and all preceding blocks.

Let $\mu_C$ be uniform probability on a cell $C$, and put
$D=a+2$, with $a$ from Lemma~\ref{lem:l-lt-likelihood}.  We will keep
the following two properties after $t$ blocks: the history fixes a
cell $C$ of depth $l$, its parameter law $P$ is regular in $C$, and
\begin{equation}\label{l:lt:density-depth}
 P\le K\mu_C,\qquad K=2^{Ddt-(p-\lambda)l}.
\end{equation}
The measure inequality means that it holds on every Borel set.
Initially $C=Q_*$, $P=\mu$, $t=l=0$, so both properties hold.

\begin{lemma}[One block and one selected cell]
\label{lem:l-lt-selection-step}
Suppose a positive-probability history after $t$ blocks has a fixed
state and a conditional law satisfying \eqref{l:lt:regular} and
\eqref{l:lt:density-depth} in a cell $C$ of depth $l$.  One can
discard conditional probability at most
\begin{equation}\label{l:lt:one-step-deletion}
 3\cdot2^{-c_1d}+(\lfloor J_t\rfloor+1)2^{-d},
 \qquad
 J_t=\frac{ad+1+Ddt}{p-\lambda},
\end{equation}
and partition the remaining outcomes into finitely many
positive-probability histories satisfying both properties at time
$t+1$.
\end{lemma}
\begin{proof}
Run the fresh next block, and let
$\mathcal E=\{q_W(z)>2^{ad}b_W\}$ be the event in
\eqref{l:lt:likelihood-tail}.  Write
$\delta=\mathbb P(\mathcal E)\le2^{-c_1d}$.  Discard $\mathcal E$.
For $b_w>0$, let
$r_w=\mathbb P(\mathcal E^c\mid W=w)$, and also discard every state
with $r_w<1/2$.  The mass of these states is at most $2\delta$:
\[
 \sum_{w:r_w<1/2}b_w
 \le2\sum_{w:r_w<1/2}\mathbb P(\mathcal E\cap\{W=w\})
 \le2\delta.
\]
The two deletions therefore cost at most $3\delta$.

For a remaining state $w$, condition on $W=w$ and $\mathcal E^c$.
Its parameter law $P'$ is given by
\[
 P'(B)=\frac1{b_wr_w}
 \int_B \mathbf1_{\{q_w(z)\le2^{ad}b_w\}}q_w(z)\,dP(z)
\]
for Borel $B$.  Since $r_w\ge1/2$,
\begin{equation}\label{l:lt:posterior-inflation}
 P'\le2^{ad+1}P\le2^{ad+1}K\mu_C.
\end{equation}

For $z\in C$, let $U_j(z)$ be the cell containing $z$ at relative
depth $j$.  Select the least $j\ge0$ maximizing the score
\begin{equation}\label{l:lt:score}
 2^{\lambda j}P'(U_j(z)).
\end{equation}
The score at $j=0$ is one.  Since
$\mu_C(U_j(z))=2^{-pj}$, \eqref{l:lt:posterior-inflation} gives
\[
 2^{\lambda j}P'(U_j(z))
 \le2^{ad+1}K2^{-(p-\lambda)j}
 \le2^{ad+1+Ddt-(p-\lambda)j}.
\]
This upper bound tends to zero.  Thus the maximum is attained, and
every selected depth is at most $J_t$, since a selected score is at
least one.  The least-maximizer rule is unambiguous.  It is Borel:
there are finitely many possible depths, the cell at each depth is a
Borel function of $z$, and each score is a fixed number on that cell.

For each possible selected cell $U$, define the actual selection event
\begin{equation}\label{l:lt:selection-event}
 F_U=\{z:\text{the rule in \eqref{l:lt:score} selects }U\}\subseteq U.
\end{equation}
These Borel events form a finite partition of $C$.  Discard those for
which
$P'(F_U)<2^{-d}P'(U)$.  At a fixed depth the cells partition $C$,
so their $P'$-masses sum to one.  Summing over the at most
$\lfloor J_t\rfloor+1$ depths bounds this deletion by the second term
of \eqref{l:lt:one-step-deletion}.  Events of zero mass need no
conditional law.

Fix a remaining event $F_U$ of positive mass, and let $j$ be the
relative depth of $U$.  Put $P''=P'(\,\cdot\mid F_U)$.  If a
descendant $V$ of $U$ at further depth $i$ meets $F_U$, choose
$z\in V\cap F_U$.  Maximality on this parameter's cell chain gives
\[
 2^{\lambda j}P'(U)\ge2^{\lambda(j+i)}P'(V).
\]
Consequently
\[
 P''(V)\le\frac{P'(V)}{P'(F_U)}
 \le2^d2^{-\lambda i},
\]
using $P'(F_U)\ge2^{-d}P'(U)$.  If $V\cap F_U$ is empty, its
$P''$-mass is zero.  This proves regularity in $U$.

The selected score is at least one, so
$P'(U)\ge2^{-\lambda j}$ and
$P'(F_U)\ge2^{-d-\lambda j}$.  The cell reference measures satisfy
$\mu_C|_U=2^{-pj}\mu_U$.  Therefore
\begin{align}
 P''
 &\le2^{ad+1}K
       \frac{2^{-pj}}{P'(F_U)}\,\mu_U\notag\\
 &\le2^{(a+1)d+1}K2^{-(p-\lambda)j}\mu_U
 \le2^{Dd}K2^{-(p-\lambda)j}\mu_U. \label{l:lt:depth-update}
\end{align}
The last inequality uses $(a+1)d+1\le(a+2)d=Dd$ for $d\ge1$.
Substituting the value of $K$ yields
\[
 P''\le2^{Dd(t+1)-(p-\lambda)(l+j)}\mu_U,
\]
which is \eqref{l:lt:density-depth} at the new time and total depth.
The normalization throughout is by $P'(F_U)$, including when
$F_U$ is a proper subset of $U$.
\end{proof}

\begin{proposition}[The finite-horizon cube success bound]
\label{prop:l-lt-forward}
For sufficiently large $d$, consider a deterministic Borel
computation padded to $m\ge0$ blocks of $n$ rows as above.  Suppose its
boundary alphabet has size at most $N$, where $N^{1/k}\le2^d$, and
its final boundary state fixes its unit output.  Its average success
under \eqref{l:lt:chart} is at most
\begin{equation}\label{l:lt:forward-bound}
 \Delta_m+2^{2Ddm}(5\epsilon)^p,\qquad
 \Delta_m=3m2^{-c_1d}
       +2^{-d}\sum_{t=0}^{m-1}(\lfloor J_t\rfloor+1).
\end{equation}
The sum is empty when $m=0$.  In particular,
$\Delta_m=o(1)$ for any sequence $m=o(d)$.
\end{proposition}
\begin{proof}
Apply Lemma~\ref{lem:l-lt-selection-step} at each surviving history
and at each block.  There are finitely many histories at every time.
The conditional deletion bounds are uniform, so their unconditional
sum is at most $\Delta_m$: at each time, the probabilities of the
surviving histories sum to at most one.  Also
$p-\lambda=(d-1)/2$, and hence $J_t=O(t+1)$ with an absolute constant
depending only on the already fixed $a$.  Thus
\[
 \Delta_m\le3m2^{-c_1d}+C(m^2+m)2^{-d},
\]
which tends to zero when $m=o(d)$.

At a surviving terminal history, integrating
\eqref{l:lt:density-depth} forces $K\ge1$.  Hence
$(p-\lambda)l\le Ddm$.  Since $\mu(C)=2^{-pl}$, the bound relative
to the original cube probability becomes
\begin{equation}\label{l:lt:terminal-density}
 P\le2^{Ddm+\lambda l}\mu\le2^{2Ddm}\mu.
\end{equation}
The last inequality uses $\lambda=p-\lambda$.  The history fixes the
final boundary state and therefore its output.  Equation
\eqref{l:lt:output-mass} bounds conditional success by
$2^{2Ddm}(5\epsilon)^p$.  Average over surviving histories and add
the deleted mass to obtain \eqref{l:lt:forward-bound}.  For $m=0$,
the same argument is the fixed-output bound at the root.
\end{proof}

The finite-horizon estimate \eqref{l:lt:forward-bound} holds for every
$m$.  We will obtain $m=o(d)$ only on the hypothetical short-run
sequence in the proof of Theorem~\ref{thm:l-lt-precision}.

\subsection{The constant-accuracy endpoint and the precision contradiction}
\label{subsec:l-lt-endgame}

The forward estimate pays a fixed amount at the first block.  A
separate full-data argument supplies the sample bound at constant
accuracy.  Even after all labels are revealed, the signal has an
undetermined direction on a residual sphere.  Its directional second
moment will bound the success probability.

\begin{lemma}[A residual-sphere sample bound]
\label{lem:l-lt-residual}
For $0<\epsilon\le1/10$ and all sufficiently large $d$, any learner
with success at least $2/3$ for every unit signal must have $T>d/4$.
This holds even if the learner is given all $T$ observations before
producing its output.
\end{lemma}
\begin{proof}
Let \(S\) be uniform on \(S^{d-1}\), independently of the \(T\)
pregenerated Gaussian rows.  The original jointly measurable experiment
has prior success equal to the integral of its fixed-signal successes.
For each admissible fixed shared seed, Lemmas~\ref{l:prior-ac} and
\ref{l:borel} give Borel rules with the same uniform-prior law.
The bound below is uniform over these rules, so it bounds the original
conditional success before averaging over the seed.  This returns to
the original learner; it does not reuse the cube-selected program.

If $T\le d/4$, the span $U$ of the rows has dimension $T$ almost
surely.  Given the rows, the labels
determine $P_US$, since the row map is injective on $U$.  Conditional
on the rows and $P_US$, the residual $(I-P_U)S$ is uniform on the
sphere in $U^\perp$ of radius
$r=(1-\|P_US\|^2)^{1/2}$.  This follows, for example, by representing
$S$ as a normalized standard Gaussian and using rotational invariance
of its component in $U^\perp$.

Spherical symmetry gives $\mathbb E\|P_US\|^2=T/d$, so Markov's
inequality yields
\[
 \mathbb P\{r<1/2\}
 =\mathbb P\{\|P_US\|^2>3/4\}
 \le\frac{4T}{3d}\le\frac13.
\]
Suppose $r\ge1/2$ and condition on the rows and labels.  Fix any
possible output, and let $w$ be its projection onto $U^\perp$.
If $w=0$, its residual error is $r>1/10$.  If $w\ne0$,
write $\Theta$ for the unit residual direction and $e=w/\|w\|$.
When $|\langle\Theta,e\rangle|<1/2$, the component of $r\Theta$
orthogonal to $e$ has length greater than
$\sqrt3\,r/2>1/10$, so the output cannot succeed.  The second moment
of $\langle\Theta,e\rangle$ for the uniform direction in the
$(d-T)$-dimensional space $U^\perp$ is $1/(d-T)$.  Markov's inequality
therefore bounds success for the fixed output by $4/(d-T)$.  The
learner's conditional output law uses only the revealed data and
independent randomness, so averaging over that law preserves the bound.
Angular success implies the Euclidean error bound used here.

The uniform-prior success is consequently at most
$1/3+4/(d-T)$, which is below $2/3$ for large $d$ when $T\le d/4$.
The every-signal premise gives uniform-prior average at least $2/3$,
a contradiction.
\end{proof}

\begin{proof}[Proof of Theorem~\ref{thm:l-lt-precision}]
Write $L=\log_2(1/\epsilon)$ and keep the absolute constants $a,c_1,D$
already fixed.  Choose
\begin{equation}\label{l:lt:endgame-constants}
 L_0\ge\max\{32D,\,2\log_2 5\},\qquad
 0<c\le\min\left\{1,\frac1{4L_0},\frac1{2048D}\right\}.
\end{equation}
Suppose for contradiction that there are infinitely many dimensions
with successful learners satisfying
\begin{equation}\label{l:lt:short-run}
 T<c\,dL.
\end{equation}
All asymptotic reductions below concern these dimensions only.

We first obtain the precision range forced by the finite set of
terminal labels.  For fixed shared rules, there are at most
$(T+1)2^M$ terminal state/index labels.  Let $\mathcal O_{t,u}$ be
the output law at one such label.  For each parameter, its success
probability is at most the sum, over all terminal labels, of the
success probabilities under their output laws: this drops the
probability of reaching each label, which is at most one.
Integrating over $z$ and then over each $\mathcal O_{t,u}$, using
\eqref{l:lt:output-mass}, bounds cube-prior success by
$(T+1)2^M(5\epsilon)^p$.  This bound is uniform in the fixed shared
rules, so it also bounds the original randomized learner.  Its
every-signal guarantee gives cube-prior average at least $2/3$.
Thus
\begin{equation}\label{l:lt:output-count}
 p(L-\log_2 5)
 \le M+\log_2(T+1)+\log_2(3/2).
\end{equation}

Because $L\ge\log_2 10$, the absolute number
$c_0=1-\log_2 5/\log_2 10$ is positive and
$L-\log_2 5\ge c_0L$.  Under \eqref{l:lt:short-run} and $c\le1$,
\[
 \log_2(T+1)\le1+\log_2 d+\log_2 L\le1+\log_2 d+L.
\]
For large $d$, absorb the last $L$ into $c_0pL$ on the left of
\eqref{l:lt:output-count}.  It follows that
\begin{equation}\label{l:lt:precision-reduction}
 L=O\!\left(\frac{M+\log d}{d}\right)=o(d),
 \qquad T=o(d^2).
\end{equation}
The constants in this bound are absolute.  In particular, for a
fixed memory sequence it bounds every $L$ arising in
\eqref{l:lt:short-run} by the same sequence $o(d)$.

At each such dimension, Lemma~\ref{lem:l-lt-borel} supplies a
deterministic Borel learner with cube-prior success at least $3/5$.
Pad a halted path by retaining its terminal state and its terminal
index.  At any later layer, the continuing states together with all
terminal tags have total width at most
\begin{equation}\label{l:lt:global-width}
 N=(T+2)2^M.
\end{equation}
Pad the last block with ignored rows and put $m=\lceil T/n\rceil$.
By \eqref{l:lt:precision-reduction},
$\log_2 N=o(d^2)$ and $m=o(d)$.  Since $k=\Theta(d)$,
$N^{1/k}\le2^d$ for all sufficiently large dimensions under
consideration.  Proposition~\ref{prop:l-lt-forward} applies and gives
\begin{equation}\label{l:lt:asymptotic-bound}
 \frac35\le \Delta_m+2^{2Ddm}(5\epsilon)^p
            =o(1)+2^{2Ddm}(5\epsilon)^p.
\end{equation}
The $o(1)$ is uniform over the short-run precisions for a fixed memory
sequence, by the uniform bound on $L$ just obtained.

If $L\le L_0$, then \eqref{l:lt:short-run} and
$c\le1/(4L_0)$ give $T<d/4$.  Apply
Lemma~\ref{lem:l-lt-residual} to the original learner and its
uniform-sphere average, obtaining a contradiction for large $d$.
This use of the original learner preserves the prior of that lemma.

It remains to consider $L>L_0$.  For large $d$, $n\ge d/64$, so
\[
 m\le64cL+1,\qquad
 2Ddm\le128Dc\,dL+2Dd\le dL/8.
\]
The last inequality follows from $c\le1/(2048D)$ and
$L\ge32D$.  Also $p=d-1\ge d/2$ and $L\ge2\log_2 5$, whence
\[
 p(L-\log_2 5)\ge dL/4.
\]
Therefore the second term in \eqref{l:lt:asymptotic-bound} is at most
$2^{-dL/8}$.  It tends to zero, as does $\Delta_m$, contradicting
the fixed lower bound $3/5$.  No infinite subsequence
\eqref{l:lt:short-run} exists.  This proves the asserted sample
bound with the absolute $c$ in \eqref{l:lt:endgame-constants}.
\end{proof}
